\begin{proposition}[D.4.7; Proved]
The function $D_{\rm part}$ has no critical point on the punctured
root neighbourhood $0<|\tau-i|\leq3/20000$, no critical point on the
coarse neighbourhood represented by $|\tau-\rho|\leq1/50$, and no
critical point where $\Im\tau\geq2$. At $\rho$ the two factors
$Z$ and $F=3ZP-2g_2$ of D.4.4 vanish, but the coarse gradient is
nonzero, as in D.4.5. The assertions away from $i$ are certified by
the independently implemented Arb calculation
\path{code/d_root_gradient_certificate.py}; all its intervals and
the 400 local rectangles are retained.
\end{proposition}
\begin{proof}
We first treat $i$ by analytic estimates. Use the holomorphic
functions $G(h,i+x)$ and $H(h,i+x)$ of E.4.2. That proof gives
$|G|<20$, $|H|<0.6$ on the polydisc with radii $1/5$ in $h$ and
$1/10$ in $x$, and gives all its roots over $|x|\leq3/20000$ in
$|h|^2<1/20000$. Put
\[
 a=\frac3{2000},\qquad b=\frac1{800},\qquad
 U=(1-a)^{-2}(1-b)^{-1}-1,\qquad
 V=(1-a)^{-1}(1-b)^{-2}-1.
\]
Differentiate the convergent even power series of $G$. Subtracting
the known coefficients $G_x(0,i)=id$ and
$G_{h^2}(0,i)=-g_2(i)/4$, Cauchy's coefficient inequalities give
\[
 |G_x-id|\leq200U<0.853,\qquad
 \frac{|G_h|}{|h|}\leq2(47.30+500V)<98.611.
\]
The first bound follows by summing
$\sum_{j\geq1,k\geq0}j a^{j-1}b^k$ after its constant term
has been removed. The second uses
$\sum_{j\geq0,k\geq1}k a^j b^{k-1}$ with the same removal.
Since $d>10.03$, $G_x$ is nonzero and the zero curve is locally
$x=x(h)$, with
\[
 |x'(h)|\leq K|h|,\qquad
 K=\frac{2(47.30+500V)}{10.03-200U}<10.745<11.
\]
At $h=0$, E.4.2 forces $x=0$; elsewhere $h$ is thus a valid local
coordinate of this curve. For
$R=H-H(0,i)+\pi h^2/2$, its constant and both linear coefficients
in $(x,h^2)$ vanish. The same sums imply
\[
 |R_h|\leq30V|h|,\qquad |R_x|\leq6U.
\]
Write $T=1+\Im x$ and $y=\Im h$. The norm on the zero curve is
$\log\mathcal N=\Re H-\pi y^2/T$. Its Wirtinger derivative satisfies
\[
 \partial_h\log\mathcal N
 =-\frac\pi2\bar h+\mathcal E,
\]
\[
 \frac{|\mathcal E|}{|h|}
 \leq\frac{30V+6UK}{2}
 +\frac{\pi\epsilon}{1-\epsilon}
 +\frac{\pi K}{40000(1-\epsilon)^2}
 <0.198862<\frac\pi2,\qquad \epsilon=\frac3{20000}.
\]
Indeed the last two terms bound respectively
$\pi i y(T^{-1}-1)$ and
$-\pi i y^2x'(h)/(2T^2)$. This proves strict nonvanishing for
$h\ne0$. Since $D=(6/\pi)\log(\mathcal N/(2\sqrt3))$, it proves
the assertion at $i$ without differentiating a numerical upper bound.

We give the details of the other two certificates. Write $q=e^{2\pi
i\tau}$ and use $\wp,\wp'$ as in E.4.1, together with
\[
\begin{aligned}
 L&=\pi\cot\pi z+
       4\pi\sum_{n\geq1}\frac{q^n}{1-q^n}\sin(2\pi nz),\\
 E&=\frac{\pi^2}{3}-8\pi^2\sum_{n\geq1}\frac{nq^n}{1-q^n},&
 g_2&=\frac{4\pi^4}{3}
       +320\pi^4\sum_{n\geq1}\frac{n^3q^n}{1-q^n}.
\end{aligned}
\]
These follow directly by differentiating the Jacobi product and the
Eisenstein series. The code sums 40 terms. For $0\leq r<1$, set
\[
 V_j(r)=41^j r^{41}M_j(r),\qquad
 (M_0,M_1,M_2,M_3)=
 \left(\frac1{1-r},\frac1{(1-r)^2},
 \frac{1+r}{(1-r)^3},\frac{1+4r+r^2}{(1-r)^4}\right).
\]
Then $\sum_{n>40}n^j r^n\leq V_j(r)$, by
$(40+k)^j\leq41^jk^j$. If $r\geq|q|$ and
$R=r\exp(2\pi|\Im z|)<1$, the tails for
$(\wp,\wp',L,E,g_2)$ are enclosed by complex squares whose
coordinate radii are respectively
\[
 \frac1{1-r}\bigl(
 8\pi^2(V_1(r)+V_1(R)),\ 16\pi^3V_2(R),\
 4\pi V_0(R),\ 8\pi^2V_1(r),\ 320\pi^4V_3(r)\bigr).
\]
They enclose the corresponding complex discs and so are conservative.
Every use of a strict inequality is an Arb comparison between
outward-rounded real balls at 256-bit precision.

For the cusp, let
\[
 z_\infty=\frac12-\frac{i}{2\pi}\log(5+2\sqrt6).
\]
The $q=0$ trigonometric function has a simple zero there. Let $z_0$
be its stored exact dyadic approximation and let $A$ be the stored
dyadic approximation of the inverse derivative at $q=0$. The
closed square $|\Re(z-z_0)|,|\Im(z-z_0)|\leq1/1000$ is invariant
under $z\mapsto z-A\wp(z,q)$ for every $q$ in the square
$|\Re q|,|\Im q|\leq e^{-4\pi}$. The certified supremum-norm
Lipschitz bound is less than $0.024725$ and the image radius is less
than $0.000122623$. Banach's theorem gives the unique analytic root
there, agreeing with the chosen cusp branch. The square in $q$
contains the entire required disc; its larger absolute-value bound
is used in every tail. In the additional interval
$0\leq1/\Im\tau\leq1/2$, direct evaluation gives
\[
 |Z|>1.39894,\qquad |3ZP-2g_2|>127.018.
\]
Consequently neither critical factor vanishes anywhere above height
two. The opposite zero $-z$ changes the sign of $Z$ and $P$ and
leaves $F$ unchanged; translations by periods leave the same test
invariant. Thus the choice of the zero representative does not
restrict this conclusion.

At $\rho$, put $z_\rho=(2-\rho)/3$ and $P_0=\wp'_\rho(z_\rho)$.
The square $\rho+[-1/50,1/50]+i[-1/50,1/50]$ is divided into
$20\times20$ closed rectangles of side $1/500$. This is an exact
rational partition relative to $\rho$. On each rectangle the code
constructs an exact dyadic centre and preconditioner and proves the
same strict contraction criterion for a root square of radius
$1/1000$. Every resulting root lies in the common closed square of
coordinate radius $1/80$ about the dyadic centre of $z_\rho$; the
largest required coordinate radius is less than $0.007731<1/80$.
On this larger square the same uniform derivative calculation
gives $\|1-A\wp'\|_\infty<0.432523<1$ for every modulus in the
whole square about $\rho$. This proves injectivity of $\wp$ there:
if two values agree, the difference of the inputs equals their
difference under $z\mapsto z-A\wp(z)$, contradicting the strict
Lipschitz inequality. This larger square need not be invariant
under that map, and we do not assert that it is. Existence on each
small box and uniqueness in the common square therefore glue one
holomorphic root. The certified nonzero values of $P$ also exclude
a double root. At the centre this is $z_\rho$.

All needed total derivatives on $\wp(\mu)=0$ can be evaluated
without numerical differentiation. With $b=\Im\mu/\Im\tau$,
they are
\[
\begin{aligned}
 \mu'&=\frac{iL}{2\pi}-\frac{i g_2}{6\pi P},\\
 Z_\tau&=\frac{-P-2EL}{4\pi i}-E\mu'
                  +\frac\pi{\Im\tau}(\mu'-b),&
 Z_{\bar\tau}&=\frac\pi{\Im\tau}(b-\overline{\mu'}),\\
 P_\tau&=\frac{-3EP-Lg_2/2}{2\pi i}-\frac{g_2}2\mu',&
 g_2'&=\frac{2iEg_2+3iP^2}{\pi},\\
 F_\tau&=3Z_\tau P+3ZP_\tau-2g_2',&
 F_{\bar\tau}&=3Z_{\bar\tau}P.
\end{aligned}
\]
The first two follow from D.4.4 and the theta heat equation. For the
third, differentiate its displayed partial-$\tau$ identity for
$\wp$ in $\mu$, then add $\wp''\mu'=-g_2\mu'/2$.
The fourth is Ramanujan's identity together with $P^2=-g_3$ at a
zero of $\wp$. In particular these formulas are identities for
the chosen analytic root, not independent-variable derivatives.

Let $A_0=iP_0/(4\pi)$ and $B_0=-21iP_0^2/(4\pi)$. The
400 strict enclosures give, throughout the square,
\[
\begin{aligned}
 |Z_\tau-A_0|+|Z_{\bar\tau}|&<0.423916
                         <2.279897<|A_0|,\\
 |F_\tau-B_0|+|F_{\bar\tau}|&<337.953
                         <1371.701<|B_0|.
\end{aligned}
\]
Since $Z(\rho)=F(\rho)=0$ by D.4.5, integrate each derivative
along the segment from $\rho$ to $\tau$. Its difference from
$A_0(\tau-\rho)$, respectively $B_0(\tau-\rho)$, has strictly
smaller absolute value than that linear term. Hence both factors
are nonzero for $0<|\tau-\rho|\leq1/50$. At the centre itself,
D.4.5 gives the nonzero coarse derivative $7/(64\pi)$.
This completes all local and cusp assertions.
\end{proof}
